\documentclass[11pt,a4paper]{article}
\usepackage[T1]{fontenc}\usepackage[utf8]{inputenc}\usepackage[brazil]{babel}
\usepackage{lmodern}\usepackage{microtype}\usepackage{amsmath,amssymb}
\usepackage{booktabs}
% max width so encolhe quando a tabela nao cabe; as estreitas ficam intactas.
\usepackage{adjustbox}
\usepackage{array}\usepackage{siunitx}\usepackage{pgfplotstable}
\usepackage[hidelinks]{hyperref}\usepackage[margin=2.7cm]{geometry}
\pgfplotsset{compat=1.18}
\sisetup{output-decimal-marker={,}, group-separator={.}}
\newcommand{\dat}[1]{dados/#1}
% termos que o babel brasileiro nao hifeniza sozinho
\hyphenation{Dis-cre-ti-za-tion HiG-Flow MF-Sim par-ti-cio-na-men-to}

\title{Extens\~ao CPU/GPU do HFlow:\\
       as premissas de desempenho, confrontadas com a medida}
\author{Relat\'orio t\'ecnico}
\date{Setembro de 2026}

% Sem isto o TeX prefere estourar a margem a afrouxar a linha, e nao
% emite aviso alto para que o log acuse.  Com emergencystretch ele usa
% os pontos de quebra dos identificadores longos em \texttt.
\emergencystretch=3em
\begin{document}
\maketitle

\begin{abstract}
\sloppy
O projeto \emph{Extens\~ao da Arquitetura HFlow para Execu\c{c}\~ao Acelerada em
CPU/GPU em Malhas Hier\'arquicas} prop\~oe tr\^es coisas: um n\'ucleo de software
compartilhado entre HiGFlow e MFSim-NG, o t8code como op\c{c}\~ao de
particionamento, e execu\c{c}\~ao em GPU no \emph{Solver} e no operador mais
custoso de \emph{Discretization}.

Este documento confronta as \emph{premissas de desempenho} desse projeto com
medidas feitas no c\'odigo. \textbf{N\~ao \'e uma avalia\c{c}\~ao do projeto},
cujo n\'ucleo \'e arquitetural e independe do que segue: unificar duas
implementa\c{c}\~oes divergentes tem valor pr\'oprio. \'E a contribui\c{c}\~ao de
uma medida para a \emph{ordem} das extens\~oes.

Tr\^es achados. O remalhamento custa entre \SI{0.06}{\percent} e
\SI{0.10}{\percent} do tempo, de modo que a troca do mecanismo de
particionamento otimiza, na escala medida, algo negligenci\'avel. A premissa de
que o ajuste de m\'inimos quadrados do GFD \'e o operador mais caro
\textbf{nunca foi medida}, e n\~ao h\'a instrumento no c\'odigo que separe
montagem de resolu\c{c}\~ao. E os gargalos estruturais que existem est\~ao do lado
da interface que o \emph{princ\'ipio de n\~ao invasividade} promete n\~ao tocar
--- uma tens\~ao a decidir, n\~ao um defeito a corrigir.
\end{abstract}

\tableofcontents
\section{As medidas que esta an\'alise usa}
\label{sec:fase0}

Duas campanhas de medida no HiGFlow~\cite{sousa2019} sustentam o que segue. Ambas foram feitas
sobre o caso de bolha ascendente bif\'asica (\emph{benchmark} de Hysing~\cite{hysing2009}), com
cerca de 9\,400 c\'elulas --- e essa pequenez \'e a principal limita\c{c}\~ao
deste documento, dita aqui e retomada no fim.

\subsection{Custo do remalhamento e escalabilidade}

O crit\'erio de refino e a reconstru\c{c}\~ao do dom\'inio foram cronometrados
\textbf{separados do resto}, com barreira antes de cada marca, em
$n_p = 1$, 2, 4 e 8:

{\small
\pgfplotstabletypeset[
  string type,
  columns/np/.style={column name={$n_p$}},
  columns/parede/.style={column name={parede}},
  columns/aceleracao/.style={column name={acelera\c{c}\~ao}},
  columns/eficiencia/.style={column name={efici\^encia}},
  columns/criterio/.style={column name={crit\'erio}},
  columns/reconstrucao/.style={column name={reconstr.}},
  columns/remalhatotal/.style={column name={remalha/total}},
  every head row/.style={before row=\toprule, after row=\midrule},
  every last row/.style={after row=\bottomrule},
]{\dat{fase0-escalonamento.dat}}
}

\subsection{Itera\c{c}\~oes do solver linear}

\pgfplotstabletypeset[
  string type,
  columns/np/.style={column name={$n_p$}},
  columns/pressao/.style={column name={press\~ao}},
  columns/vsbase/.style={column name={vs.\ base}},
  columns/velocidade/.style={column name={velocidade}},
  every head row/.style={before row=\toprule, after row=\midrule},
  every last row/.style={after row=\bottomrule},
]{\dat{fase0-ksp.dat}}

A press\~ao domina o custo linear com folga --- 110 itera\c{c}\~oes contra 4 a 6
da velocidade ---, como se espera da equa\c{c}\~ao de Poisson da proje\c{c}\~ao.
O crescimento de \SI{11}{\percent} com o n\'umero de processos \'e a
degrada\c{c}\~ao do precondicionador \texttt{bjacobi} ao ser dividido em mais
blocos, e responde por cerca de \SI{13}{\percent} da acelera\c{c}\~ao perdida.


H\'a um projeto de pesquisa em prepara\c{c}\~ao --- \emph{Extens\~ao da
Arquitetura HFlow para Execu\c{c}\~ao Acelerada em CPU/GPU em Malhas
Hier\'arquicas} --- que prop\~oe tr\^es coisas: um n\'ucleo de software
compartilhado entre HiGFlow e MFSim-NG, o t8code como nova op\c{c}\~ao de
particionamento, e uma camada de execu\c{c}\~ao em GPU no \emph{Solver} e no
operador mais custoso de \emph{Discretization}.

Este documento confronta as premissas de desempenho desse projeto com o que
foi medido no c\'odigo. \textbf{N\~ao \'e uma avalia\c{c}\~ao do projeto}, cujo
n\'ucleo \'e arquitetural e independe do que segue; \'e a contribui\c{c}\~ao de
uma medida para a \emph{ordem} das extens\~oes.

\section{As tr\^es premissas de desempenho}

{\small
\pgfplotstabletypeset[
  string type,
  columns/premissa/.style={column name={premissa do projeto},
    column type={>{\raggedright\arraybackslash}p{5.0cm}}},
  columns/estado/.style={column name={estado},
    column type={>{\centering\arraybackslash}p{1.9cm}}},
  columns/medida/.style={column name={o que a medida diz},
    column type={>{\raggedright\arraybackslash}p{6.0cm}}},
  every head row/.style={before row=\toprule, after row=\midrule},
  every last row/.style={after row=\bottomrule},
]{\dat{projeto-premissas.dat}}
}

\section{O que foi medido}

\pgfplotstabletypeset[
  string type,
  columns/grandeza/.style={column name={grandeza}},
  columns/valor/.style={column name={valor}},
  columns/onde/.style={column name={procedência}},
  every head row/.style={before row=\toprule, after row=\midrule},
  every last row/.style={after row=\bottomrule},
]{\dat{projeto-medido.dat}}

\section{Tr\^es consequ\^encias}

\paragraph{1. O t8code otimizaria, hoje, algo que custa \SI{0.1}{\percent}.}
Isso \textbf{n\~ao refuta} a extens\~ao: o projeto mira malhas muito maiores que
a medida aqui, e nesse regime o equil\'ibrio muda. O que a medida estabelece \'e
que o ganho \'e \emph{invis\'ivel na escala atual}, e portanto que a
justificativa n\~ao pode ser verificada sem antes levar um caso ao tamanho-alvo.

\paragraph{2. A premissa da GPU n\~ao est\'a medida, e a \'unica evid\^encia
dispon\'ivel aponta para o outro lado.} O projeto identifica o ajuste de m\'inimos
quadrados do GFD como o operador mais caro. Montagem e resolu\c{c}\~ao
\textbf{nunca foram cronometradas separadamente} --- n\~ao h\'a instrumento no
c\'odigo que as separe. O que h\'a \'e a contagem de itera\c{c}\~oes: a press\~ao
consome 110 itera\c{c}\~oes de KSP por passo contra 4 a 6 da velocidade, o que
sugere peso na \emph{resolu\c{c}\~ao}.

Pode ser que o MLS domine mesmo assim --- montagem cara e poucas
itera\c{c}\~oes n\~ao se excluem. O ponto \'e que a prioriza\c{c}\~ao entre
\emph{Solver} e \emph{Discretization} hoje repousa sobre asser\c{c}\~ao, e que
separ\'a-las custa um cron\^ometro.

\paragraph{3. E aqui est\'a o ponto mais delicado.} Os tr\^es gargalos
estruturais identificados na leitura do c\'odigo s\~ao:

\begin{enumerate}
  \item uma \'unica \'arvore \texttt{.amr} materializada \emph{inteira num
    processo} antes do particionamento;
  \item o dom\'inio \textbf{reconstru\'ido do zero} a cada remalhamento, em vez de
    adaptado no lugar;
  \item o crit\'erio de refino passando por \textbf{arquivo escrito
    serialmente}, lido por todos.
\end{enumerate}

Nenhum deles est\'a no \emph{backend} de particionamento. Est\~ao na camada da
aplica\c{c}\~ao e do n\'ucleo --- do lado da interface que o \emph{princ\'ipio de
n\~ao invasividade} do projeto promete n\~ao tocar.

Trocar MTree por t8code atr\'as da mesma interface de consulta topol\'ogica
melhora a qualidade da parti\c{c}\~ao de uma malha que continua sendo montada
globalmente num processo. \textbf{O crit\'erio de projeto que protege a
arquitetura \'e, nesse sentido, o que impede a extens\~ao de alcan\c{c}ar os
gargalos} --- e isso \'e uma tens\~ao real a decidir, n\~ao um defeito a corrigir:
n\~ao invasividade e realoca\c{c}\~ao de gargalo s\~ao objetivos que competem.

\section{O que isto sugere para a ordem do trabalho}

O pr\'oprio projeto prev\^e, na fase F1, um \emph{levantamento de
\emph{baseline}: \emph{benchmarks} de desempenho e de escalabilidade da
arquitetura atual (CPU, MTree)}. A sugest\~ao que decorre desta se\c{c}\~ao \'e
trat\'a-lo como \textbf{decis\'orio, e n\~ao preliminar}: fazer a
\emph{baseline} no tamanho-alvo antes de fixar a ordem das extens\~oes, porque
ela pode relocar o gargalo.

Tr\^es medidas resolveriam as d\'uvidas acima, e as tr\^es s\~ao baratas:

\begin{itemize}
  \item \textbf{montagem $\times$ resolu\c{c}\~ao}, cronometradas em separado ---
    decide se a GPU come\c{c}a pelo \emph{Solver} ou pelo \emph{Discretization};
  \item \textbf{escalabilidade forte num caso grande} --- o medido aqui tem cerca
    de 9\,400 c\'elulas, e a \SI{22}{\percent} de efici\^encia a conclus\~ao \'e
    sobre o \emph{tamanho}, n\~ao sobre a arquitetura;
  \item \textbf{mem\'oria por processo ao aumentar $n_p$} --- \'e o teste direto do
    primeiro gargalo, e o \'unico que diria se a materializa\c{c}\~ao global morde
    antes do resto.
\end{itemize}

\section{O alcance desta se\c{c}\~ao, e os seus limites}

Convém delimitar, porque a se\c{c}\~ao faz afirma\c{c}\~oes sobre um projeto a
partir de medidas estreitas:

\begin{itemize}
  \item \textbf{um \'unico caso, e pequeno}: bolha ascendente bif\'asica, cerca de
    9\,400 c\'elulas. Nada aqui se estende automaticamente a malhas grandes, que
    s\~ao o alvo do projeto;
  \item \textbf{um \'unico software}: tudo foi medido no HiGFlow. O MFSim-NG usa
    FV e n\~ao GFD, e a premissa sobre o operador mais caro pode ser verdadeira
    l\'a e falsa aqui, ou o contr\'ario;
  \item \textbf{a arquitetura n\~ao foi avaliada}. O n\'ucleo compartilhado \'e a
    parte principal do projeto, e nada nesta se\c{c}\~ao o toca --- unificar duas
    implementa\c{c}\~oes divergentes tem valor independente de desempenho;
  \item \textbf{os tr\^es gargalos s\~ao fatos de c\'odigo; que eles
    \emph{dominem} \'e infer\^encia}, e a medida da invent\'ario do sistema
    mostrou que infer\^encias desse tipo j\'a caíram uma vez neste trabalho.
\end{itemize}

\bibliographystyle{plain}
\bibliography{referencias}
\end{document}
